\documentclass[10pt,a4paper]{amsart}
\usepackage[margin=19mm]{geometry}
\usepackage{amsmath,amssymb,mathtools}
\usepackage[scr=rsfs,cal=euler]{mathalfa}
\usepackage{xfrac}
\usepackage[unicode,hidelinks,bookmarks=false]{hyperref}
\newcommand{\ie}{i.e.\ }
\newcommand{\cf}{cf.\ }
\newcommand{\resp}{respectively\ }
\newcommand{\Subsection}{\subsection}
\providecommand{\nobreakdash}{\nobreak\hbox{-}}
\setlength{\emergencystretch}{2em}
\multlinegap0pt
\usepackage{bm}
\usepackage{bbm}
\usepackage{dsfont}
\usepackage{xspace}
\usepackage{cancel}
\usepackage{mathtools}
\usepackage{amsthm}
\newtheorem{thm}{Theorem}
\title[Generalized Goulden-Yong duals and signed minimal factorizat]{Generalized Goulden-Yong duals and signed minimal factorizations}
\author{Shujian Chen, Kiyoshi Igusa}
\date{Source: arXiv:2311.05732v3 (CC BY 4.0)}
\begin{document}
\maketitle

\noindent\emph{Source and licence.} Generalized Goulden-Yong duals and signed minimal factorizations. Version arXiv:2311.05732v3 (CC BY 4.0), distributed under the Creative Commons Attribution 4.0 International licence, as recorded in the arXiv licence metadata for this exact version and shown on the abstract page https://arxiv.org/abs/2311.05732v3. Full rights notice: licenses/arxiv-2311.05732-CC-BY-4.0.txt.\par\smallskip

\noindent\textit{Editorial scope.} Counted unit: the Thm whose source label is \texttt{Bernardi} in the article (source0.tex lines 1657--1660, source label \texttt{Bernardi}), delivered together with the article's own complete original proof of it (source0.tex lines 1662--1694, one \texttt{proof} environment of 33 lines). No mathematics is written, repaired or re-derived: statement and proof are the article's own text line for line. Required context retained uncounted: none. Every definition, display and label of those lines is retained unchanged. Omitted from this package and not redistributed: 1--1656, 1661--1661, 1695--2083. Nothing inside a retained excerpt is omitted or altered: every retained body is delivered exactly as the article's lines are written, so no omission receipt of any kind accompanies this package. Citations: the retained excerpts carry no citation command at all, so no bibliography entry of the article is transcribed and no reference is redistributed. Reference adaptation: none was needed and none was made; every reference inside the delivered unit resolves inside it, so no unresolved reference or citation is produced. Printed slips kept literal: no printed word, symbol, hypothesis or number of the counted statement or of its proof was altered. Layout and class: the article's own class is \texttt{amsart} with packages including geometry, graphicx, amsthm, amsfonts, tikz, indentfirst, amsmath, amssymb; the wrapper is the lane's pinned \texttt{amsart} editorial wrapper at 10pt on A4, which restates the theorem-like environments the retained text needs and adds the article's own macro definitions that the retained text needs (none), with extra wrapper packages (bm, bbm, dsfont, xspace, cancel, mathtools). The delivered numbering is the unit's own. Assets: no retained span contains a figure environment, a graphics-inclusion command, a PDF-page inclusion, a table or a TikZ picture; no image file is copied into this package, the package declares no asset dependency and no image file is redistributed with it. Licence: version arXiv:2311.05732v3 of this article is distributed under the Creative Commons Attribution 4.0 International licence, as recorded in the arXiv licence metadata for this exact version and shown on the abstract page https://arxiv.org/abs/2311.05732v3. A full rights notice is delivered at licenses/arxiv-2311.05732-CC-BY-4.0.txt. Characters: every retained excerpt is pure ASCII, so the delivered engine, XeLaTeX with its default Latin Modern text font, reports no missing character.
\begin{thm}[Bernardi] \label{Bernardi}
    Let $S=\{$rooted trees on $n$ vertices such that the root is smaller than its children$\}$ and 
    $T=\{$rooted trees on $n$ vertices such that $v_n$ is a leaf$\}$. There's a bijection between $g: S \rightarrow T$.
\end{thm}

\begin{proof}
    We'll first describe the map $g$:
    
    Take $s \in S$, let $v_k$ be the root of $s$.

    Take the rooted forest of all descendants of $v_n$ where the roots are the children of $v_n$ and move it under $v_k$, i.e., make descendants of $v_n$ become descendants of $v_k$ while preserve the tree structure.

    There's a unique path $P$ between $v_k$ and $v_n$ in $s$.

    For all children of $v_k$ that are not on the path, we can arrange them in descending order, i.e. there's a sequence $(v_{k_1},\dots,v_{k_j})$ such that $k_i > k_{i+1}$.

    Make $v_{k_{i+1}}$ a child of $v_{k_i}$ for all $1 \leq i \leq j$ while preserving all other descendants.
    
    Make $v_{k_1}$ the root of $s$ instead of $v_k$. This completes the construction of $g(s)$.\\

    Now we will describe $g^{-1}$:

    Take $t \in T$, let $v_k$ be the root of $t$.

    There's a unique path $P$ between $v_k$ and $v_n$ in $t$.

    Order the vertices on the path $P$ as $(v_k,v_{i_1},v_{i_2},\dots,v_n)$ such that $v_{i_j}$ is a child of $v_{i_{j-1}}$.

    There exists an unique element $v_{i_l}$ such that $v_k > v_{i_1} > \dots > v_{i_l}$. (If the whole sequence is ascending, take $v_{i_l}$ to be $v_k$)

    Take the rooted forest of all descendants of $v_{i_l}$ where the roots are the children of $v_{i_l}$ and move it under $v_n$, i.e., make descendants of $v_{i_l}$ become descendants of $v_n$ while preserve the tree structure.

    For elements on the left of $v_{i_l}$ in the sequence, i.e. $(v_k,v_{i_1},\dots,v_{i_{l-1}})$, make each of them a child of $v_{i_l}$ along with all their descendants.

    Now $v_{i_l}$ is the root for the tree.

    This complete the construction of $g^{-1}$.
\end{proof}

\end{document}
